%% frontmatter/preface.tex
%% Preface — explicit four-section structure: who / how / prereqs / outcomes

\chapter*{Preface}
\addcontentsline{toc}{chapter}{Preface}

\PLACEHOLDER{The preface is a tour for the reader, not the author.
A short version follows; the full text will be written once the chapters are filled in.}

\section*{Who this book is for}

This book is written for the curious reader with a basic feel for probability and algebra, who wants to understand why \emph{random matrices} are everywhere in modern data science --- from principal component analysis to deep learning, from finance to wireless communications --- and who has been told the words ``Tracy--Widom,'' ``spiked covariance,'' or ``spectral clustering'' without ever being shown how they fit together.

If you have seen a normal distribution and you know what an eigenvalue is, you have enough background to start. Everything else is built up from scratch.

\section*{How to read this book}

The book supports two reading speeds.

\begin{description}
  \item[The \emph{cliff-notes track}.] Read the chapter opener (\textsc{Where we are} box), the \textsc{Intuition} boxes, the figures, and the \textsc{Chapter recap} at the end. Skip the proofs. You will come away with the conceptual picture in roughly 30\% of the time.
  \item[The \emph{deep-diver track}.] Read everything: definitions, theorems, proofs, examples, simulations, exercises. The book is structured so that every claim is either proved in-line, sketched with a pointer to a full proof in the chapter appendix, or quoted as a standard result with a citation.
\end{description}

The right margin is reserved for asides, small figures, code snippets, and Colab badges. \emph{Asides never carry plot-load} --- if you ignore the entire margin, the proofs still work. Read the margins for context, history, motivation, or the occasional small joke.

\PLACEHOLDER{Cliff-notes track listing: explicit list of which sections to read. To be written once chapters are filled in.}

\section*{What you need to know first}

\begin{itemize}
  \item Some idea of what a random variable is, and what a probability density looks like (we'll re-build all of this from scratch in Part~I, but having seen it once before will speed things up).
  \item Comfort with linear algebra at the level of vectors, matrices, eigenvalues, and a dot product. We re-derive the spectral theorem from scratch in Chapter~6.
  \item Basic Python literacy --- enough to read a NumPy snippet and run a Colab notebook.
\end{itemize}

\section*{What you'll be able to do after}

After reading this book, you will:

\begin{itemize}
  \item Read research papers on random matrix theory and follow the central results without getting stuck on definitions.
  \item Run and interpret simulations of sample covariance matrices, eigenvalue distributions, and Marchenko--Pastur and Tracy--Widom limits.
  \item Recognize the spiked covariance model and the BBP phase transition in real data.
  \item Implement spectral clustering and explain \emph{three} different reasons it works.
\end{itemize}

\section*{A note on style}

The book is written in third person, in the style of polished research notes that have been opened up to a wider audience. Technical content is rigorous; prose is friendly. Cartoon asides exist. The title is a small joke; the math is not.

\bigskip
\noindent\hfill\textit{Arzaan Singh}\par
\noindent\hfill New Orleans, \DTMtoday
